\section{Discussion}
\label{sec:discussion}
Houdini’s design not only addresses critical security challenges in container confinement but also lends itself well to continuous security testing in modern software development environments. Because it is both safe and modular, \houdini can be seamlessly integrated into automated systems, such as CI/CD pipelines, to provide ongoing validation of container security.

Regression testing is a widely adopted practice in software engineering to ensure that updates or changes to code do not introduce unintended failures or security regressions. In the context of container security, integrating \houdini into CI/CD pipelines enables organizations to automatically trigger a comprehensive suite of confinement tests whenever any part of the container stack changes—such as kernel upgrades, Docker releases, or policy updates. This continuous validation ensures that every new deployment or system update is accompanied by an up-to-date security assessment. By systematically rerunning empirical security tests across different software versions and configurations, \houdini helps detect configuration drift and regression in security guarantees that might otherwise go unnoticed, maintaining the effectiveness of container isolation mechanisms as platforms evolve.